
\section{Contraste bibliográfico: fuerza hidrostática y centro de presiones}
\label{sec:ref26-hidro}

La presión de un fluido homogéneo en reposo, medida respecto de la presión que actúa al otro lado de la superficie, es \(p_{\mathrm{man}}(h)=\rho gh\) para profundidad vertical \(h\) medida desde la superficie libre. Se supone densidad constante, gravedad uniforme y superficie libre con presión de referencia uniforme. No debe confundirse la presión manométrica con la absoluta: si ambas caras de una compuerta están sometidas a la misma presión atmosférica, esa contribución se cancela en la fuerza neta.

En una placa \emph{vertical}, rectangular, de anchura \(b\), altura \(H\) y borde superior a profundidad \(h_0\), la fuerza y su punto de aplicación se obtienen integrando la distribución, no evaluando la presión máxima:
\begin{equation}
 F=\int_0^H \rho g(h_0+y)b\,dy
   =\rho g A h_c,\quad
 A=bH,\quad h_c=h_0+\frac{H}{2}.
 \label{eq:ref26-hidro-fuerza}
\end{equation}
Al igualar momentos respecto a la superficie libre,
\begin{equation}
 h_{\mathrm{cp}}
 =\frac{\int_A h^2\,dA}{\int_A h\,dA}
 =h_c+\frac{I_G}{Ah_c},\qquad I_G=\frac{bH^3}{12},
 \label{eq:ref26-hidro-cp}
\end{equation}
donde \(I_G\) es el segundo momento del área respecto del eje \emph{horizontal centroidal}. La segunda expresión presupone una superficie plana vertical; para una placa inclinada debe introducirse el ángulo geométrico correspondiente. El centro de presiones queda por debajo del centroide si la superficie está totalmente sumergida.

\begin{tcolorbox}[title={Ejemplo guiado original: compuerta}, studiumdef, breakable]
Una compuerta vertical de \(1{,}20\,\mathrm{m}\) de anchura y \(1{,}50\,\mathrm{m}\) de altura tiene el borde superior a \(0{,}60\,\mathrm{m}\) de profundidad. Para agua con \(\rho=1000\,\mathrm{kg\,m^{-3}}\) y \(g=9{,}81\,\mathrm{m\,s^{-2}}\), se obtiene \(A=1{,}80\,\mathrm{m^2}\), \(h_c=1{,}35\,\mathrm{m}\) y \(I_G=0{,}3375\,\mathrm{m^4}\).
La resultante vale \(F=1000\cdot9{,}81\cdot1{,}80\cdot1{,}35=\mathbf{23838{,}3\,\mathrm{N}}\).
La profundidad de su línea de acción es
\[
 h_{\mathrm{cp}}=1{,}35+\frac{0{,}3375}{1{,}80\cdot1{,}35}
 =\mathbf{1{,}48889\,\mathrm{m}},
\]
es decir, \(0{,}88889\,\mathrm{m}\) por debajo del borde superior. Se verifica \(h_c<h_{\mathrm{cp}}<h_0+H\), como exige la distribución triangular/trapezoidal de presiones. Las unidades de \(I_G/(Ah_c)\) son metros.
\end{tcolorbox}

\textbf{Autocomprobación.} ¿Qué ocurre si el borde superior coincide con la superficie libre (\(h_0=0\))? Demuestre que \(h_{\mathrm{cp}}=2H/3\) y no \(H/2\). Explique por qué una misma presión atmosférica sobre ambas caras no cambia el resultado. La deducción es propia; el tratamiento general de las distribuciones de presión se contrasta con White, capítulo 2 \cite{ref26-white}.
